\chapter{Asbestos Management Risk Assessment}
\label{chap:Asbestos Management Risk Assessment}

%---------------------------- SECTION ----------------------------
\section{Why this assessment matters in construction}

Asbestos is the single largest cause of occupational cancer mortality in the United Kingdom. Approximately 5,000 people die from asbestos-related diseases every year in the UK --- more than the total number of deaths on UK roads --- and the death toll is projected to continue at this level until the 2030s as the legacy of widespread asbestos use in the post-war construction boom works through its 30 to 40-year latency period. Mesothelioma, lung cancer, asbestosis, and pleural thickening are all caused by the inhalation of asbestos fibres; all are largely incurable; and all are caused, in the context of construction, by the disturbance of asbestos-containing materials (ACMs) during maintenance, refurbishment, and demolition work on buildings constructed before the mid-1980s.

The tragedy of asbestos-related disease in the construction industry is that it is entirely preventable. Unlike many occupational health hazards, the mechanism of harm (fibre inhalation), the materials that contain the hazard (ACMs in pre-2000 buildings), and the activities that cause exposure (drilling, cutting, grinding, or otherwise disturbing ACMs) are all well understood. The Control of Asbestos Regulations 2012 provide a comprehensive and clear regulatory framework for managing asbestos in construction. The HSE's guidance publications --- most importantly L143 (Managing and Working with Asbestos) --- set out in precise detail what is required. And yet, the HSE continues to identify asbestos management failures as one of the most common serious non-compliances it encounters during inspection of refurbishment and maintenance contractors.

The persistence of asbestos deaths and enforcement failures reflects several specific challenges. The first is the presence of asbestos in a vast number of buildings that their current owners, managers, and occupants may not know about: buildings constructed between 1930 and 1980 commonly contain asbestos in textured coatings (artex), floor tiles, ceiling tiles, pipe lagging, insulating board, roof sheets, and countless other applications. The second is the existence of what the HSE calls the ``cowboy'' refurbishment sector: small contractors who do not survey for asbestos before commencing alteration work, do not recognise ACMs when they encounter them, and disturb them without protection or notification to the HSE. The third is the failure of the duty to manage asbestos in non-domestic premises: many duty holders do not have an up-to-date asbestos management survey and management plan, without which contractors working on the premises cannot know what ACMs are present or where they are.

\barquote{``Every tradespeople who drilled or cut their way through textured coatings and ceiling tiles in the 1970s and 1980s without protection is now in the latency period of a disease they cannot avoid. The question now is whether we repeat that mistake with the next generation.''}

\example{Site Reality}{A small electrical contractor is engaged to install a new consumer unit and rewire a 1970s-built commercial property. The building does not have an asbestos management survey. The electrician, believing the suspended ceiling tiles to be plasterboard, drills 47 fixing points through ceiling tiles to run cables in the ceiling void. The tiles are Asbestolux insulating board containing between 15 and 25 per cent chrysotile asbestos. The electrician has no respiratory protection, no enclosure, and no decontamination procedure. Airborne fibre counts taken by an independent analyst two days later, during a repeat attendance, record levels significantly above the clearance indicator. The electrician is 34 years old; the consequences of the exposure will not be known for 30 years. The electrical contractor's RAMS made no reference to asbestos surveys or pre-survey confirmation; the building owner had no asbestos management survey and had not told the contractor about the building's age or construction.}

%---------------------------- SECTION ----------------------------
\section{Legal and professional framework}

The Control of Asbestos Regulations 2012 (CAR 2012) are the primary legislative instrument for asbestos management in construction. They impose three fundamental duties: the duty to manage asbestos in non-domestic premises (Regulation 4); the duty to identify the type and condition of asbestos before any work that may disturb it is undertaken; and a licensing regime for the most hazardous types of asbestos work. CAR 2012 establishes three categories of asbestos work: licensed work (highest-risk, requires HSE licence, notification to HSE, Health and Safety Plan, and specific worker health surveillance); notifiable non-licensed work (NNLW, which must be notified to the relevant enforcing authority and requires health surveillance); and non-licensed work (the lowest-risk work on asbestos, which still requires COSHH assessment, RPE, and safe working procedure).

HSE publication L143 (Managing and Working with Asbestos) provides the comprehensive compliance guidance for CAR 2012. The type of asbestos survey required is determined by the purpose of the work: a management survey (type 1) is required for the ongoing management of ACMs in occupied premises; a refurbishment and demolition survey (type 2) is required before any refurbishment, alteration, or demolition work that may disturb the building fabric.

\paragraph{Key frameworks relevant to this chapter include: Control of Asbestos Regulations 2012; HSE L143; COSHH Regulations 2002; CDM 2015.}

\begin{itemize}
  \bitem{Control of Asbestos Regulations 2012}{Impose the duty to manage asbestos in non-domestic premises; require identification of ACMs before work that may disturb them; establish the licensed, NNLW, and non-licensed work categories with associated requirements; require health surveillance for all workers who carry out any work with asbestos; prohibit the importation, supply, and use of all forms of asbestos.}
  \bitem{HSE L143 Managing and Working with Asbestos}{Provides definitive compliance guidance for CAR 2012; explains the survey requirements, the licensed/NNLW/non-licensed classification framework, the health surveillance requirements, the training requirements (asbestos awareness for all who may disturb ACMs in the course of their work; CAT2 for NNLW; specific competency requirements for licensed work), and the requirements for RPE, decontamination, and waste disposal.}
  \bitem{CDM 2015}{Requires the designer to provide pre-construction information identifying the presence of known ACMs; requires the principal contractor to ensure that ACMs are identified in the Construction Phase Plan and that appropriate asbestos management arrangements are in place before any work that may disturb ACMs commences.}
  \bitem{COSHH Regulations 2002}{Apply to asbestos as a substance hazardous to health; the COSHH assessment for asbestos work must address the type and form of asbestos, the nature of the disturbance activity, the likely fibre release, and the controls required to achieve adequate control of exposure.}
\end{itemize}

%---------------------------- SECTION ----------------------------
\section{Conducting the assessment using the five-step method}

\subsection{Step 1 --- Identify hazards}
\paragraph{The primary hazard in asbestos management is the inhalation of airborne asbestos fibres released when ACMs are disturbed. ACMs likely to be encountered in UK buildings constructed before 2000 include: sprayed asbestos coatings on structural steelwork and concrete (high fibre release, typically licensed work); asbestos insulating board (AIB) used as ceiling tiles, partition panels, fire breaks, and soffits (high fibre release, typically NNLW or licensed); asbestos cement in roof sheets, gutters, downpipes, and external panels (moderate fibre release, typically non-licensed); pipe and boiler lagging in thermal insulation (variable, depending on fibre type and condition); floor tiles (chrysotile in vinyl floor tiles, low fibre release unless tiles are friable or activities generate significant dust); textured decorative coatings (artex, typically chrysotile, low to moderate fibre release when sanded or drilled); gaskets and rope seals in plant and pipework systems. The hazard is amplified by the type of fibre (crocidolite/amosite more dangerous than chrysotile), the condition of the ACM (damaged, friable ACMs release more fibres), and the nature of the activity.}

\subsection{Step 2 --- Decide who or what is affected}
\paragraph{Persons at risk from asbestos exposure in construction include: directly exposed tradespeople (electricians, plumbers, plasterers, carpenters, roofers) working on pre-2000 buildings without prior survey confirmation; asbestos removal operatives (licensed and NNLW); building occupants if ACMs are disturbed without adequate enclosure and fibre control; adjacent building occupants if removal enclosures are not adequately contained; waste transfer personnel handling asbestos waste in skips or bags without appropriate containment. Risk to the general public arises where asbestos removal is conducted at open construction sites without adequate hoarding and air monitoring.}

\subsection{Step 3 --- Evaluate likelihood and consequence}
\paragraph{The consequence of significant asbestos fibre inhalation is potentially fatal (mesothelioma, lung cancer), with a latency period of 20 to 50 years. This characteristic makes the risk unique among construction hazards: the affected worker will not know the outcome of an exposure incident for decades, and the disease, once manifested, is not reversible or curable. The inherent risk for undiscovered or uncontrolled ACM disturbance must be rated at the maximum (Very High, 25/25); even non-licensed work, which involves lower fibre release activities, requires proper assessment and controls to maintain residual risk at Low.}

\subsection{Step 4 --- Select and implement controls}
\paragraph{The primary control for asbestos in construction is the pre-work refurbishment and demolition survey: commission an UKAS-accredited surveyor to conduct a type R\&D survey of the entire area to be worked in before any work that may disturb the building fabric commences. The survey identifies all ACMs, their type, condition, and extent. Licensed ACMs must be removed by a licensed contractor before other work commences. NNLW operations must be conducted by trained operatives with health surveillance, correct RPE (minimum FFP3), and appropriate decontamination procedures. Non-licensed work must be subject to a COSHH assessment. All asbestos waste must be double-bagged in UN-approved packaging, clearly labelled, and disposed of at a licensed waste carrier and tip.}

\subsection{Step 5 --- Record and review}
\paragraph{The asbestos survey must be retained and made available to all contractors working on the premises; a copy must be incorporated into the Construction Phase Health and Safety File. Health surveillance records for all workers who have conducted licensed or NNLW asbestos work must be retained for 40 years. Licensed work notifications to the HSE must be retained. Air monitoring records during and after licensed removal works must be retained. The asbestos register must be updated whenever an R\&D survey identifies new materials or whenever the condition of ACMs changes.}

\example{Five-Step Application}{A refurbishment contractor is undertaking the full strip-out of a 1960s-built school before conversion to offices. Step 1 identifies: AIB ceiling tiles in classrooms; sprayed asbestos coating on steelwork in the roof void; asbestos cement on external rainwater goods. Step 2 identifies: all trades working in the building; licensed asbestos removal team; air monitoring surveyor. Step 3 rates undiscovered ACM disturbance risk as Very High (25/25). Step 4 implements: full UKAS-accredited R\&D survey completed before any strip-out work commences; licensed contractor engaged for sprayed coating removal and AIB removal; notification submitted to HSE; licensed removal completed and air clearance certificate obtained before any other trade enters those areas; non-licensed asbestos cement work conducted by trained operatives with health surveillance and FFP3 RPE. Residual risk rated Low (4/25) after licensed removal. Step 5: clearance indicator air test and certificate retained in H\&S file.}

%---------------------------- SECTION ----------------------------
\section{Applying the hierarchy of controls}

\paragraph{The hierarchy for asbestos management in construction is largely determined by CAR 2012: the primary control is to identify all ACMs before work commences and to remove high-risk ACMs by licensed contractors before other works proceed. Controls for residual exposures then follow the general hierarchy.}

\begin{itemize}
  \bitem{Elimination}{Commission a refurbishment and demolition survey before any work commences; remove all ACMs found to be in poor condition or at risk of disturbance by a licensed contractor before other work begins; design new construction so that no future maintenance or refurbishment work will expose workers to ACMs.}
  \bitem{Substitution}{For non-licensed asbestos cement work (roof sheets, drainage components), assess whether replacement with non-ACM alternatives is preferable to work on the existing ACMs; replacement eliminates the residual management obligation.}
  \bitem{Engineering controls}{Fully enclosed asbestos removal enclosure with negative pressure air unit and continuous air monitoring during licensed removal; local exhaust ventilation (LEV) or vacuum system with HEPA filtration for NNLW operations; wet suppression techniques to reduce fibre release during non-licensed work.}
  \bitem{Administrative controls}{Asbestos register and management plan for the premises maintained and provided to all contractors; mandatory R\&D survey before any refurbishment or demolition; competency requirements (awareness training for all who may encounter ACMs; CAT2 for NNLW; licensed operative competency for licensed work); health surveillance for all licensed and NNLW workers; notification to HSE for licensed work; UN-approved asbestos waste packaging and disposal.}
  \bitem{PPE}{Disposable overalls (Type 5/6) for all asbestos work; FFP3 or higher RPE (half-mask P3 or powered air purifying respirator for NNLW; supplied-air breathing apparatus for high-risk licensed work); gloves; boot covers; decontamination procedure (shower unit for licensed work) before leaving the work area.}
\end{itemize}

%---------------------------- SECTION ----------------------------
\section{Cadence of assessment and review triggers}

\paragraph{The asbestos management process has a defined lifecycle driven by the type of work and the survey findings; review triggers are specific and well defined in CAR 2012 and L143.}

\begin{itemize}
  \bitem{Before every refurbishment or demolition project}{A type R\&D survey must be commissioned and completed before any work commences; this is not negotiable or subject to programme pressure; no work that may disturb the building fabric may commence until the survey is complete and reviewed.}
  \bitem{When survey findings change}{If additional ACMs are discovered during work (discovery of unlisted materials), work must stop, the material must be sampled by an accredited analyst, and the risk assessment must be updated before work resumes in the affected area.}
  \bitem{When the ACM condition changes}{If a management survey identifies a change in the condition of an ACM (deterioration from intact to damaged), the risk assessment must be updated and the management plan revised to address the increased risk.}
  \bitem{At the commencement of each licence period}{Asbestos licensed contractor licences are renewed every three years; the licence holder's competency, training records, health surveillance, and safety management system must be reviewed at each renewal.}
  \bitem{Following any uncontrolled exposure event}{Any incident in which workers may have been exposed to asbestos fibres without adequate controls must trigger an immediate investigation, an HSE notification, an air quality assessment of the affected area, and a review of the risk assessment.}
\end{itemize}

%---------------------------- SECTION ----------------------------
\section{Common failure modes}

\begin{itemize}
  \bitem{No pre-work survey commissioned}{The most fundamental and most common failure: refurbishment work commences without an R\&D survey; tradespeople disturb ACMs without awareness, protective equipment, or decontamination. This is the primary cause of uncontrolled asbestos exposure in the UK construction industry.}
  \bitem{Management survey used for refurbishment work}{A management survey (type 1, for ongoing management of ACMs in occupied premises) is mistakenly relied upon as equivalent to an R\&D survey for refurbishment work; management surveys are not designed to locate all ACMs, including those in locations that would only be accessible during demolition.}
  \bitem{Unlicensed contractor undertaking licensed work}{Small contractors classify work that requires a licence as NNLW or non-licensed to avoid the cost and complexity of licensed contractor engagement; the classification of work must be made by reference to CAR 2012 and L143, not by reference to commercial convenience.}
  \bitem{RPE not fitted or not appropriate}{FFP3 disposable respirators are specified but fitted incorrectly; facial fit has not been tested; pre-existing beard growth prevents adequate face-piece seal. RPE is the last line of defence in asbestos management; its adequacy must be verified by fit testing.}
  \bitem{Asbestos waste not correctly packaged or disposed of}{Asbestos debris collected in ordinary polythene bags, placed in general waste skips, or transported by general waste contractors who do not hold the necessary waste carrier registration. Asbestos waste requires double UN-approved packaging, labelling, a consignment note, and disposal at a licensed facility.}
  \bitem{Health surveillance not implemented}{Operatives conducting licensed and NNLW asbestos work are not enrolled in the health surveillance programme required by CAR 2012 and records are not retained. Health surveillance for asbestos workers must be maintained for 40 years from the last date of exposure.}
\end{itemize}

\example{Failure Pattern}{A property developer is carrying out a full refurbishment of a 1975-built factory unit. Asbestos cement roof sheets are present, correctly identified in a management survey. The developer engages a general building contractor for the internal strip-out, who removes the existing suspended ceiling AIB tiles without checking the management survey (which records the tiles as ``potential, not assessed for refurbishment''). The management survey was not a type R\&D survey. The tiles contain 18 per cent amosite asbestos. Eight operatives conduct the strip-out over three days without any RPE, enclosure, or decontamination. The situation is discovered when the asbestos analyst attending to certify asbestos cement disposal notices the AIB debris. An emergency investigation follows; eight operatives have been exposed to amosite asbestos fibres over three days.}

%---------------------------- CHAPTER SUMMARY ----------------------------
\section{Chapter Summary}

\begin{table}[H]
\centering
\begin{tabularx}{\textwidth}{>{\raggedright\arraybackslash}p{3.2cm}X}
\toprule
\textbf{Assessment Element} & \textbf{Operational Requirement} \\
\midrule
Legal/Professional Basis & Control of Asbestos Regulations 2012; HSE L143; CDM 2015. R\&D survey required before all refurbishment work; licensed contractor for licensed work; health surveillance records retained 40 years. \\
Method & Apply the full five-step process and document inherent/residual risk. \\
Controls & UKAS-accredited R\&D survey before work; licensed removal of high-risk ACMs; NNLW/non-licensed procedures for residual work; FFP3 RPE fit-tested; correct waste management. \\
Cadence & Before every refurbishment project; when new ACMs are discovered; when ACM condition changes; following any uncontrolled exposure. \\
Assurance & Use audit, incident learning, and governance reporting to verify effectiveness. \\
\bottomrule
\end{tabularx}
\end{table}

\begin{itemize}
  \bitem{Key takeaway 1}{A type R\&D asbestos survey by an UKAS-accredited surveyor must be completed before any refurbishment or demolition work that may disturb the building fabric; no other control is effective if ACMs are not first identified.}
  \bitem{Key takeaway 2}{The classification of asbestos work as licensed, NNLW, or non-licensed must be made by reference to CAR 2012 and L143, not by commercial convenience; unlicensed contractors who conduct licensed work are committing a serious criminal offence.}
  \bitem{Key takeaway 3}{Health surveillance for all licensed and NNLW asbestos workers is a legal requirement; records must be retained for 40 years from last exposure; operatives must have access to their own health surveillance records on request.}
  \bitem{Key takeaway 4}{Asbestos waste must be double-packaged in UN-approved containers, labelled with the asbestos warning symbol, transported by a registered waste carrier, and disposed of at a licensed facility; general waste disposal is illegal and dangerous.}
\end{itemize}

\barquote{In Chapter~\ref{chap:Silica, Wood Dust and Respiratory Hazard Risk Assessment}, we turn from this assessment to the next domain in the Prudentia framework.}
